% !TEX root = ../main.tex
%************************************************
\chapter{RoleBox Platform Design Specifications}
\labch{rolebox-design}
%************************************************

\begin{abstract}
This appendix documents the design requirements, technical specifications, and user experience guidelines for the RoleBox multimodal analysis platform. RoleBox provides interactive browser-based tools for exploring, coding, and visualizing organizational role data across eight modalities: websites, social media, Crunchbase, visual materials, news media, Wikipedia, role simulation, and cross-modal schema integration.
\end{abstract}

\section{Design System Foundation}

The RoleBox platform adopts IBM's Carbon Design System\sidecite{ibm-carbon-2024} as its foundational design language, ensuring visual consistency and accessibility across all tools.

\subsection{Typography}

All RoleBox interfaces use IBM Plex, an open-source typeface family designed for optimal screen legibility:

\begin{itemize}
    \item \textbf{IBM Plex Sans}: Primary typeface for UI elements, headings, and body text
    \item \textbf{IBM Plex Mono}: Monospace typeface for code, identifiers, and technical labels
\end{itemize}

\begin{marginfigure}
\begin{lstlisting}[language=CSS,basicstyle=\tiny\ttfamily]
:root {
  --font-sans: 'IBM Plex Sans',
    -apple-system, sans-serif;
  --font-mono: 'IBM Plex Mono',
    monospace;
}
\end{lstlisting}
\caption{CSS custom property definitions for RoleBox typography.}
\end{marginfigure}

\subsection{Color Palette}

The platform employs a semantic color system derived from Carbon's palette, with specific assignments for the EIT Knowledge Innovation Communities (KICs):

\begin{table}[h]
\centering
\small
\begin{tabular}{lll}
\toprule
\textbf{Variable} & \textbf{Hex Code} & \textbf{Application} \\
\midrule
\texttt{--blue-60} & \texttt{\#0f62fe} & Primary actions, links \\
\texttt{--purple-60} & \texttt{\#8a3ffc} & Crunchbase, ventures \\
\texttt{--green-60} & \texttt{\#198038} & Websites, success states \\
\texttt{--teal-60} & \texttt{\#007d79} & Social media, RoleSim \\
\texttt{--magenta-60} & \texttt{\#d02670} & Visual materials \\
\texttt{--orange-50} & \texttt{\#eb6200} & News media \\
\texttt{--cyan-50} & \texttt{\#1192e8} & Twitter/social \\
\bottomrule
\end{tabular}
\caption{RoleBox semantic color assignments by modality.}
\labtab{rolebox-colors}
\end{table}

\subsubsection{KIC Sector Colors}

Each EIT Knowledge Innovation Community receives a distinct color for data visualization consistency:

\begin{table}[h]
\centering
\small
\begin{tabular}{ll}
\toprule
\textbf{KIC Sector} & \textbf{Color} \\
\midrule
Climate-KIC & Blue (\texttt{\#0f62fe}) \\
EIT Digital & Purple (\texttt{\#8a3ffc}) \\
EIT Health & Red (\texttt{\#da1e28}) \\
EIT InnoEnergy & Yellow (\texttt{\#f1c21b}) \\
EIT Food & Green (\texttt{\#198038}) \\
EIT Manufacturing & Gray (\texttt{\#525252}) \\
EIT RawMaterials & Brown (\texttt{\#8a3800}) \\
EIT Urban Mobility & Teal (\texttt{\#007d79}) \\
\bottomrule
\end{tabular}
\caption{Color assignments for EIT KIC sectors.}
\labtab{kic-colors}
\end{table}

\subsection{Navigation Architecture}

The platform implements a hierarchical navigation structure with three levels:

\begin{enumerate}
    \item \textbf{Master Hub} (\texttt{data/index.html}): Central landing page linking all modalities
    \item \textbf{Modality Index Pages}: Hub pages for each data source (e.g., \texttt{rolebox-websites/ui/index.html})
    \item \textbf{Individual Tools}: Specialized analysis interfaces within each modality
\end{enumerate}

All pages feature breadcrumb navigation allowing users to traverse upward through the hierarchy. Tool pages include a ``back arrow'' link to their parent modality index.

%------------------------------------------------
\section{Modality-Specific Tools}
%------------------------------------------------

\subsection{RoleBox Websites}
\labsec{design-websites}

The websites modality provides tools for analyzing organizational self-presentation through website content.

\subsubsection{Design Requirements}

\begin{description}
    \item[Primary Color:] Green (\texttt{\#198038})
    \item[Chapter Reference:] Chapter 5: Interstitial Pluralism
    \item[Data Volume:] 847 sites, 10.9K hyperlinks
\end{description}

\subsubsection{Tool Suite}

\paragraph{RoleBox Triplets} validates and codes Subject-Predicate-Object triplets extracted from website text. Features include:
\begin{itemize}
    \item Human-AI collaborative coding with Model Assist integration
    \item Codebook management for predicate categories
    \item Network visualization of triplet relationships
    \item Import/export in JSON format
\end{itemize}

\paragraph{RoleBox Triads} provides per-sector Knowledge Triangle analysis using ternary plots\sidecite{powell2012-triple-helix}. Features include:
\begin{itemize}
    \item Ternary coordinate visualization (Education--Research--Business)
    \item X-Curve dynamics overlay\sidecite{hebinck2022-xcurve}
    \item Sector comparison views
    \item PNG export with transparent backgrounds
\end{itemize}

\paragraph{RoleBox Sociotypes} implements the Watch-Complexity role typology for website-based actor classification.

\paragraph{Network Layouts} visualizes inter-organizational hyperlink networks with force-directed and geographic layout options.

\subsection{RoleBox Social}
\labsec{design-social}

The social media modality analyzes Twitter/X network structures and influence patterns.

\subsubsection{Design Requirements}

\begin{description}
    \item[Primary Color:] Cyan (\texttt{\#1192e8})
    \item[Chapter Reference:] Chapter 6: Role Configurations
    \item[Data Volume:] 2,847 accounts, 156K connections
\end{description}

\subsubsection{Tool Suite}

\paragraph{RoleBox NetCodes} annotates network actors using the Watch-Complexity typology\sidecite{bruns2017-twitter-follower}. The 19-role classification includes:

\begin{margintable}
\small
\begin{tabular}{ll}
\toprule
\textbf{Category} & \textbf{Roles} \\
\midrule
High Influence & TopHub, Hub, \\
 & MegAmplifier, \\
 & Amplifier, Vulcano \\
Bridging & Bridge, Connector, \\
 & Transceiver, \\
 & Channeler, Chain \\
Flow Regulation & Dam, Reducer \\
Peripheral & Receiver, \\
 & ReceiverBranch, \\
 & BlackHole, Idle \\
Emitters & Emitter, \\
 & EmitterBranch, \\
 & LowEmitter \\
\bottomrule
\end{tabular}
\caption{Watch-Complexity role categories.}
\end{margintable}

\paragraph{RoleBox RoleField} maps actors in two-dimensional Watch-Complexity space, visualizing the distribution of network roles across the field.

\paragraph{Twitter Workflow} documents the data collection and processing pipeline for Twitter/X data.

\subsection{RoleBox Crunchbase}
\labsec{design-crunchbase}

The Crunchbase modality analyzes venture data, funding patterns, and startup ecosystem positioning.

\subsubsection{Design Requirements}

\begin{description}
    \item[Primary Color:] Purple (\texttt{\#8a3ffc})
    \item[Chapter Reference:] Chapter 7: Venture Positioning
    \item[Data Volume:] 847 ventures, \$2.4B total funding
\end{description}

\subsubsection{Tool Suite}

\paragraph{RoleBox Ventures} classifies ventures by role typology using a codebook of positioning strategies:
\begin{itemize}
    \item Technology-push vs. market-pull orientation
    \item Deep-tech vs. application focus
    \item B2B vs. B2C market positioning
    \item Climate/sustainability emphasis
\end{itemize}

\paragraph{RoleBox Network Viewer} visualizes co-investment networks, founder connections, and sector linkages using Sigma.js\sidecite{sigmajs} for WebGL-accelerated rendering.

\subsection{RoleBox Visual}
\labsec{design-visual}

The visual modality analyzes organizational imagery through semiotic register coding.

\subsubsection{Design Requirements}

\begin{description}
    \item[Primary Color:] Magenta (\texttt{\#d02670})
    \item[Chapter Reference:] Chapter 8: Visual Registers
    \item[Data Volume:] 3,456 images, 12 register types
\end{description}

\subsubsection{Tool Suite}

\paragraph{RoleBox VizReg} codes visual registers from organizational images. The coding schema captures:

\begin{itemize}
    \item \textbf{Gaze Types}: Demand, Offer, Absent, Group, Object
    \item \textbf{Setting}: Office, Lab, Outdoor, Studio, Event
    \item \textbf{Composition}: Close-up, Medium shot, Long shot
    \item \textbf{Register}: Professional, Innovation, Academic, Corporate, Startup, Community, Tech, Green
\end{itemize}

The tool supports batch processing with Ollama/LLaVA integration for AI-assisted coding.

\subsection{RoleBox News}
\labsec{design-news}

The news modality analyzes media coverage of EIT actors and innovation narratives.

\subsubsection{Design Requirements}

\begin{description}
    \item[Primary Color:] Orange (\texttt{\#eb6200})
    \item[Chapter Reference:] Appendix E: News Sources
    \item[Data Volume:] 12,456 articles, 287 sources, 8 languages
\end{description}

\subsubsection{Tool Suite}

\paragraph{News Knowledge Graph} visualizes entity co-occurrence networks extracted from news articles, supporting temporal filtering and sentiment overlay.

\subsection{RoleBox Wikipedia}
\labsec{design-wikipedia}

The Wikipedia modality provides Wikidata entity linking and knowledge graph construction.

\subsubsection{Design Requirements}

\begin{description}
    \item[Primary Color:] Blue (\texttt{\#0f62fe})
    \item[Chapter Reference:] Appendix G: Wikipedia Data
    \item[Data Volume:] 847 entities, 5.2K relations
\end{description}

\subsubsection{Tool Suite}

\paragraph{RoleBox Wikipedia KG} visualizes the knowledge graph using vis-network\sidecite{visjs}, with cluster-based navigation and entity search.

%------------------------------------------------
\section{Cross-Modal Platforms}
%------------------------------------------------

\subsection{RoleSim}
\labsec{design-rolesim}

RoleSim provides agent-based simulation of role constellation dynamics.

\subsubsection{Design Requirements}

\begin{description}
    \item[Primary Color:] Teal (\texttt{\#007d79})
    \item[Chapter Reference:] Intermezzo: Role Simulation
    \item[Capabilities:] 19 role types, configurable scenarios, temporal dynamics
\end{description}

\subsubsection{Workflow Screens}

The RoleSim interface provides eight workflow screens:

\begin{enumerate}
    \item \textbf{Dashboard}: Overview statistics and quick actions
    \item \textbf{Roles}: Watch-Complexity role definitions and visual reference
    \item \textbf{Agents}: Actor listing with role assignments
    \item \textbf{Scenarios}: Predefined simulation configurations
    \item \textbf{Simulate}: Runtime controls and progress monitoring
    \item \textbf{Results}: Role stability metrics and transition analysis
    \item \textbf{Cartography}: Spatial network visualization
    \item \textbf{Export}: Data export in multiple formats
\end{enumerate}

\subsubsection{Data Integration}

RoleSim supports multiple data sources:
\begin{itemize}
    \item Hyphe CSV import for web crawl networks
    \item Pre-generated EIT network datasets
    \item Synthetic network generation for testing
\end{itemize}

\subsection{RoleBox PyDNA}
\labsec{design-pydna}

PyDNA provides discourse network analysis and cross-modal schema integration.

\subsubsection{Design Requirements}

\begin{description}
    \item[Primary Color:] Blue (\texttt{\#0f62fe})
    \item[Reference:] Cross-Modal Schema Documentation
    \item[Data Volume:] 847 actors, 12.4K statements, 287 concepts
\end{description}

\subsubsection{Tool Suite}

\paragraph{RoleDNA} implements qualitative discourse network analysis\sidecite{leifeld2013-dna}:
\begin{itemize}
    \item Document annotation interface
    \item Statement coding (Actor--Concept--Polarity)
    \item Actor-concept network visualization
    \item Coalition detection algorithms
\end{itemize}

\paragraph{CARTALOG} provides actor catalog management:
\begin{itemize}
    \item Entity resolution across data sources
    \item Name variant tracking
    \item Cross-corpus linking (Crunchbase, Wikidata, etc.)
    \item Annotation and moderation workflow
    \item REST API for programmatic access
\end{itemize}

%------------------------------------------------
\section{Technical Specifications}
%------------------------------------------------

\subsection{Browser Requirements}

RoleBox tools are designed for modern browsers supporting:
\begin{itemize}
    \item ES6+ JavaScript (async/await, modules)
    \item CSS Custom Properties (variables)
    \item CSS Grid and Flexbox layout
    \item Canvas and WebGL (for network visualization)
\end{itemize}

Minimum supported versions: Chrome 80+, Firefox 75+, Safari 13+, Edge 80+.

\subsection{Visualization Libraries}

The platform integrates several JavaScript visualization libraries:

\begin{table}[h]
\centering
\small
\begin{tabular}{lll}
\toprule
\textbf{Library} & \textbf{Version} & \textbf{Application} \\
\midrule
Plotly.js & 2.27.0 & Charts, scatter plots \\
Sigma.js & 2.4.0 & Large network graphs \\
Graphology & 0.25.4 & Graph data structures \\
vis-network & 9.1.6 & Interactive networks \\
D3.js & 7.x & Custom visualizations \\
\bottomrule
\end{tabular}
\caption{JavaScript visualization libraries used in RoleBox.}
\labtab{viz-libraries}
\end{table}

\subsection{AI/ML Integration}

Several tools support AI-assisted coding through:
\begin{itemize}
    \item \textbf{Ollama (Local)}: Self-hosted LLM inference for privacy-sensitive data
    \item \textbf{OpenAI API}: GPT-4 integration for high-accuracy suggestions
    \item \textbf{Anthropic API}: Claude integration for nuanced classification
    \item \textbf{LLaVA}: Multimodal vision-language model for image analysis
\end{itemize}

\subsection{Data Formats}

All tools support JSON as the primary data interchange format. Export options include:
\begin{itemize}
    \item JSON (structured data)
    \item CSV (tabular data)
    \item PNG (visualizations with transparency)
    \item GraphML/GEXF (network data)
\end{itemize}

%------------------------------------------------
\section{User Experience Guidelines}
%------------------------------------------------

\subsection{Demo Data Loading}

All tools auto-load demonstration data on page initialization, ensuring users immediately see a populated interface rather than empty states. This supports:
\begin{itemize}
    \item Rapid onboarding and feature discovery
    \item Visual documentation of expected data structures
    \item Testing and development workflows
\end{itemize}

\subsection{Keyboard Shortcuts}

Tools implement consistent keyboard shortcuts where applicable:
\begin{itemize}
    \item \texttt{Ctrl+S}: Save/export current state
    \item \texttt{Ctrl+Z}: Undo last action
    \item \texttt{Arrow keys}: Navigate between items
    \item \texttt{Enter}: Confirm selection
    \item \texttt{Escape}: Close modal/panel
\end{itemize}

\subsection{Responsive Design}

While optimized for desktop use (1280px+ viewport), tools implement responsive breakpoints for tablet viewing. Mobile support is limited due to the data-intensive nature of analysis tasks.

\subsection{Accessibility}

The platform follows WCAG 2.1 AA guidelines where feasible:
\begin{itemize}
    \item Sufficient color contrast ratios (4.5:1 minimum)
    \item Keyboard navigation support
    \item Focus indicators for interactive elements
    \item Screen reader compatibility for core workflows
\end{itemize}

%------------------------------------------------
\section{Deployment Architecture}
%------------------------------------------------

RoleBox tools are implemented as static HTML/CSS/JavaScript files, deployable via:
\begin{itemize}
    \item Local file system (development/research use)
    \item Static web hosting (GitHub Pages, Netlify)
    \item Institutional web servers
\end{itemize}

The CARTALOG tool additionally requires a Python backend (FastAPI) for database operations and REST API functionality.

\begin{kaobox}[title=Platform Access]
The complete RoleBox platform is accessible from the dissertation data directory at \texttt{data/index.html}. Individual modality hubs are located at \texttt{data/rolebox-[modality]/ui/index.html}.
\end{kaobox}
